\section{Step 10: Match sequence and the full-frame schedule}
The four sequence bytes are serialized big-endian. Read them in order with
\texttt{lbu}, shifting the accumulated value left eight bits before adding the
next byte. Use 64-bit shifts so that the unsigned four-byte value remains
zero-extended. A native \texttt{lw} would depend on native byte order, require
alignment, and sign-extend some values.

Fill \texttt{record\_sequence}. Compare the result with the saved expected
sequence. Compute its unsigned remainder by the saved period, which was already
proved nonzero. A zero remainder requires kind zero. A nonzero remainder permits
either kind; full frames between scheduled positions remain legal. Sequence four
with period four therefore cannot be delta, whereas sequence five can be full.

\paragraph{Observe} Prepare \texttt{s10-1} and stop before
\texttt{record\_sequence\_ready}. At \texttt{record\_sequence\_ready}, \texttt{\$t0} is one
and the remainder in \texttt{\$t2} is one. The worked record is allowed.

\Needspace{7\baselineskip}
\paragraph{Check} With period four, which kind is valid at sequence five?
\begin{itemize}
\item[A.] Only delta.
\item[B.] Either full or delta.
\item[C.] Neither kind.
\end{itemize}
